\scanpage{88}
\mspage{8}

schéma en carrés combinatoire sur $S$) dont le schéma des diagonales s'identifie à $\underline\omega_{\underline J}$ :

\[
\underline Q_{E/S}/\mathord{\mathrm{I}\mkern-7mu\mu}_2\simeq \underline\omega_{\underline J}
\tag*{(42)}
\]

Ceci posé, on voit que l'on a un isom \inserted{can} de $(\mathord{\mathrm{I}\mkern-7mu\mu}_4)_S$-torseurs

\[
Q_{E/S}\mathbin{\wedge}_{\mathord{\mathrm{I}\mkern-7mu\mu}_4}Q_{0S}\simeq\mathcal L_{E/S}
\tag*{(43)}
\]

i.e

\[
\deleted{Q}\mathcal L_{E/S}\simeq\underline{\mathrm{Isom}}_{\mathord{\mathrm{I}\mkern-7mu\mu}_{4S}\text{-tors.}}(Q_{0S}^{\deleted{\unsure{4|\otimes|\natural}}{-1}},Q_{E/S})
\]

On aurait pu définir les épinglages de Legendre de $E/S$ comme les trivialisations de $\deleted{\mathcal L}Q_{E/S}$, au lieu de celles de $\mathcal L_{E/S}$ -- mais on aura alors une relation moins jolie avec les épinglages de Jacobi d'échelon $4$.

On voit en tous cas que les épinglages de \deleted{T}Legendre \inserted{de $E/S$} forment \inserted{les sections d'}un torseur sous $\mathord{\mathrm{I}\mkern-7mu\mu}_4$, tel\deleted{\unsure{le|s|les}} que le $\mathord{\mathrm{I}\mkern-7mu\mu}_2$-torseur associé soit can. iso. à $\tikz[baseline=(w.base)]{\node[inner sep=1pt] (w) at (0,0) {$\underline\omega_{\underline J}$};\node[inner sep=1pt] at (.68,0) {$\wedge$};\node[inner sep=1pt] (m) at (1.26,0) {$\mathord{\mathrm{I}\mkern-7mu\mu}_4^*$};\draw (w.south west) .. controls (.35,-.5) and (.65,-.36) .. (.68,0) .. controls (.68,.3) and (.95,.48) .. (m.north east);}$

\[
\mathcal L_{E/S}/\mathord{\mathrm{I}\mkern-7mu\mu}_2\simeq\mathcal L_{E/S}\mathbin{\wedge}_{\mathord{\mathrm{I}\mkern-7mu\mu}_{4\deleted{\unsure{S}}}}\mathord{\mathrm{I}\mkern-7mu\mu}_{2\deleted{S}}\simeq\deleted{\underline\omega_{\underline J}}\bigl(\underbrace{\mathord{\mathrm{I}\mkern-7mu\mu}_{4\deleted{S}}^{*}\mathbin{\wedge}\underline\omega_{\underline J}}
\tag*{(44)}
\]

tandis que les épinglages de Jacobi d'échelon

% archive sheet number: 87
